\subsection{An unconditional refinement by an admissible common source}
\label{sec:A-paths}

We use the following precise category throughout this section. For a
commutative ring $k$, put $S=\Gamma_k(T)$, with $T$ of weight one
and form degree zero. An object $A$ is a simplicial augmented strict
PD differential $S$-algebra $A^{w,q}$, with $w,q\geq0$ and differential
$d:A^{w,q}\to A^{w,q+1}$. For each $d_0\geq0$ there are integers
$\lambda_A(d_0),\mu_A(d_0)$ with the following properties:
\begin{enumerate}
\item For every $w\geq\lambda_A(d_0)$ and every simplex degree $n$,
$S^w\to H^0(A_n^{w,*})$ is an isomorphism and
$H^q(A_n^{w,*})=0$ for $1\leq q\leq d_0$.
\item For every $w\geq\mu_A(d_0)$ and $0\leq q\leq d_0$, the
nonnegative homological normalization $N(A^{w,q})$ is exact, including
homological degree zero.
\end{enumerate}
Morphisms preserve both gradings, coefficients, augmentation,
differential, multiplication and all divided powers. We denote this
category by $\mathcal A_k$. No completion is used. We also consider its
full subcategories with $A^{w,q}=0$ for $q>w$, and, over $\Z$, with
flat homogeneous terms.

\begin{lemma}[A-001: Proved; recheck of the part of SA08 C02 used here]
\label{A001}
Let $M$ be a simplicial $k$-module with exact normalization. For every
inclusion $L\hookrightarrow K$ of simplicial sets, the restriction
\[
 \Hom_{s\operatorname{Mod}_k}(k[K],M)
 \longrightarrow \Hom_{s\operatorname{Mod}_k}(k[L],M)
\]
is surjective. For every simplicial set $K$ and every $i>0$, we have
$\Ext^i_{s\operatorname{Mod}_k}(k[K],M)=0$.
Consequently $M$ is flasque in SGA~4, expos\'e V, Definition~4.1
\cite{Root_SGA4}.
\end{lemma}
\begin{proof}
We give the lifting argument without a finiteness restriction on $K$.
Under the exact Dold--Kan equivalence
\cite[Theorem 14.24.3]{Root_DoldKan}, the normalized quotient
$P=N(k[K]/k[L])$ is degreewise free, on the nondegenerate simplices of
$K$ outside $L$. A degreewise split extension
$0\to NM\to E\to P\to0$ has a graded section $s$. Its differential
defect is a degree-minus-one chain cocycle with values in $NM$.
We correct $s$ inductively, beginning in homological degree zero.
At degree $n$, the discrepancy after the earlier corrections is a
map $P_n\to Z_{n-1}(NM)$. Exactness gives a surjection
$(NM)_n\to Z_{n-1}(NM)$; freeness of $P_n$ lifts the discrepancy.
Subtracting that lift corrects the section in degree $n$ without
altering earlier degrees. Thus every extension splits as complexes.
All extensions are degreewise split because each $P_n$ is projective.
Hence $\Ext^1(P,NM)=0$. Applied to the pullback extension of a
prescribed map $k[L]\to M$, this gives the required restriction lift.

For higher Ext, resolve a degreewise-projective nonnegative complex
$P$ by identity disks. The canonical projective epimorphism is
\[
 P_0[0]\oplus\bigoplus_{n\geq1}D^n(P_n)\longrightarrow P,
\]
where $D^n(V)$ has identity differential in degrees $n,n-1$.
Its kernel has term $P_{n+1}$ in degree $n$ and differential $-d_P$,
so its terms are again projective. Repeated dimension shifting reduces
every positive Ext group to the splitting argument just proved.
This applies to $P=N(k[K])$ and to its successive kernels.

For the projection $p_K:\widehat\Delta/K\to\widehat\Delta$, the
left adjoint on module sheaves is pointwise direct sum over the fibers
and is exact. Adjunction on injective resolutions therefore gives
\[
 H^i(\widehat\Delta/K,p_K^*M)
 =\Ext^i_{s\operatorname{Mod}_k}(k[K],M)=0.
\]
This is exactly the definition in the cited source, for every $K$.
\end{proof}

\begin{proposition}[A-002: Proved; admissible path refinement]
\label{A002}
Every cospan $A\xrightarrow{f}C\xleftarrow{g}B$ in $\mathcal A_k$
has an object $H\in\mathcal A_k$ and morphisms $H\to A$, $H\to B$.
The construction preserves the support restriction $q\leq w$.
Over $k=\Z$ it also preserves termwise flatness.
\end{proposition}
\begin{proof}
For a simplicial $k$-module $M$, define its path module by
\[
 (M^I)_n=\Hom_{s\operatorname{Mod}_k}(k[\Delta[n]\times\Delta[1]],M).
\]
Endpoint restriction gives $e_0,e_1:M^I\to M$, and constant paths give
$c:M\to M^I$. These maps are defined in each weight and form degree.
Pointwise products and divided powers make $C^I$ a strict PD DGA.
The augmentation is well defined because $\Delta[n]\times\Delta[1]$
is connected, so its maps to a constant simplicial module are constant.

Write $P_0M=\ker(e_0:M^I\to M)$. The simplicial contraction of the
interval to its initial vertex contracts $P_0M$ to zero. Explicitly the
order-preserving minimum map $[1]\times[1]\to[1]$ defines a simplicial
homotopy fixing the initial vertex; precomposition gives this additive
contraction. No flatness is needed.

We also need the vertical comparison. Fix $d_0$ and take $w$ beyond
the finitely many resolution and row bounds for $C$ through $d_0+3$.
The coefficient map $S^w\to C^{w,*}$ is then a resolution through
that degree on the simplicial topos, whose relevant terms are flasque
by Lemma~\ref{A001}. The two first-quadrant hypercohomology spectral
sequences identify its section cohomology through degree $d_0+1$
with constant-coefficient cohomology on any simplicial set. Only those
finitely many form degrees contribute to this range; the bound does
not depend on the dimension or cardinality of the simplicial set.
Apply this to $\Delta[n]\times\Delta[1]$ and its initial endpoint.
They are contractible and the endpoint inclusion is a homotopy
equivalence. Restriction is termwise surjective in the relevant range
by Lemma~\ref{A001}. Its kernel $P_0C$ is therefore vertically exact
through degree $d_0$, uniformly in $n$.

Form the mapping-path object and its endpoint map
\[
 X=A\mathop{\times}_{C,e_0}C^I,
 \qquad q:X\to C,\quad q(a,\alpha)=e_1\alpha.
\]
The assignment $(a,\alpha)\mapsto(a,\alpha-cf(a))$ is an isomorphism
of simplicial bigraded $k$-modules $X\cong A\oplus P_0C$.
Its inverse is $(a,\beta)\mapsto(a,\beta+cf(a))$.
We use this isomorphism only additively. The preceding horizontal
contraction and eventual vertical exactness prove $X\in\mathcal A_k$.
In the eventual vertical range, the map $q$ induces the identity on
the specified coefficient $S^w$ and is an isomorphism on all cohomology.

Set $H=X\times_C B$ in simplicial PD DGAs and let $K=\ker q$.
In any fixed form range and sufficiently large weights, endpoint
restriction $C^I\to C\times C$ is termwise surjective by
Lemma~\ref{A001}. Thus $q:X\to C$ and $H\to B$ are termwise
surjective in that range. We obtain short exact sequences of the
relevant simplicial rows and of the relevant vertical truncations:
\[
 0\to K\to X\xrightarrow{q}C\to0,
 \qquad 0\to K\to H\to B\to0.
\]
All three rows $X,C,B$ are exact in high weights. The first long exact
homology sequence makes $K$ exact, and the second makes $H$ exact.
Vertically the first sequence makes $K$ exact in the fixed range
because $q$ is the coefficient-preserving cohomology isomorphism
just proved. The second identifies the cohomology of $H$ with that
of $B$ in that range. Increasing the chosen bound by one form degree
whenever a connecting map requires it proves both admissibility
conditions, uniformly in $n$. Projection through $X\to A$ gives
the two claimed strict morphisms. All objects and maps retain both
gradings; vanishing above $q=w$ is therefore preserved.

Finally work over $\Z$ and assume flat terms. The normalized integral
chains of the finite contractible set $\Delta[n]\times\Delta[1]$
split as $\Z[0]$ and finitely many free identity disks: every cycle
subgroup is free, and the exact cycle sequences split. Chain maps
from these summands into $NC^{w,q}$ form finite sums of normalized
terms and cycles. These are torsion free because the terms of $C$
are torsion free. Hence every term of $C^I$, then of $X$ and $H$,
is torsion free, as a subgroup of a finite product of torsion-free
groups. A torsion-free abelian group is flat. This argument covers
all weights, including those below the eventual bounds.
\end{proof}

\begin{corollary}[A-003: Proved; precise relation between common sources and connectedness]
\label{A003}
In $\mathcal A_k$, two objects belong to the same connected component
of the nerve exactly when they have an admissible common source.
This also holds with the support condition, and in the termwise-flat
subcategory over $\Z$. Therefore terminality of the localization at
all morphisms entails common-source existence in these categories.
\end{corollary}
\begin{proof}
A common source supplies a two-arrow zigzag. Conversely, connectivity
is represented by a finite zigzag. Inductively retain a span from
one object to its current endpoint and to its initial endpoint.
A forward arrow composes with the first leg. For a backward arrow,
apply Proposition~\ref{A002} to the cospan formed by that leg and
the backward arrow, and compose its new projection with the retained
leg to the initial endpoint. This constructs the required span after
each step. All compositions are morphisms in the indicated category.
A localization at all arrows has the same connected components as
these finite zigzags. A terminal localization is connected.
\end{proof}

The corollary proves the stated implication by explicit morphisms.
It supplies neither an initial zigzag between arbitrary integral
objects nor a contraction of a zigzag loop. Those are the two tasks
still requiring construction in Part A.
